\begin{abstract}
We prove that every finite loopless \(k\)-edge-connected multigraph on
at least two vertices, with \(k\ge1\), has a spanning tree meeting each
cut in at most \(C/k\) times the size of the cut,
where \(C\) is a universal constant.
This resolves the strong thin tree conjecture.
\end{abstract}

\section{Introduction}\label{sec:introduction}

For a finite loopless undirected multigraph \(G=(V,E)\) and a nonempty proper
subset \(S\subset V\), write \(\delta_G(S)\) for the edges with exactly
one endpoint in \(S\). Parallel edges are counted separately.
A graph is \(k\)-edge-connected if \(|\delta_G(S)|\ge k\) for
every nonempty proper \(S\).
A spanning tree \(T\) is \(\alpha\)-thin if
\[
 |\delta_T(S)|\le\alpha|\delta_G(S)|
 \qquad(\varnothing\ne S\subsetneq V).
\]
A spanning tree retains connectivity using only \(|V|-1\) edges.
Thinness asks that these edges occupy a small fraction of every cut
simultaneously. The question is whether high edge connectivity alone
forces such a tree, with a bound independent of the number of vertices.

\begin{theorem}\label{thm:main}
There is a universal constant \(C>0\) such that, for every integer
\(k\ge1\), every finite \(k\)-edge-connected undirected multigraph
without loops and with at least two vertices has a spanning tree \(T\)
satisfying
\[
 |\delta_T(S)|\le\frac Ck\,|\delta_G(S)|
 \qquad\text{for every }\varnothing\ne S\subsetneq V.
\]
\end{theorem}
